\documentclass[12pt]{article}
\usepackage[utf8]{inputenc}
\usepackage[T1]{fontenc}
\usepackage{hyperref}
\title{Turkish Sign Language Language Resources for K-12 Curriculum Development}
\author{Compiled by Claude Sonnet 5 (claude-sonnet-5) for Andr\'as Kornai}
\date{}
\begin{document}
\maketitle

\section{Introduction}
Turkish Sign Language (T\.{I}D; ISO 639-3 tsm, Glottocode turk1288) is the sign language of Turkey's deaf community. This report is based on a single search, so coverage is limited.

\section{Primary Educational Resources}
\begin{itemize}
\item Bogazici University work on web-based material to improve Turkish reading and writing for T\.{I}D users (a paper in Turkish): \url{https://www.cmpe.boun.edu.tr/~gungort/papers/Turk%20Isaret%20Dili%20(TID)%20Kullanicilarinin%20Turkce%20Okuma%20ve%20Yazma%20Becerilerinin%20Gelistirilmesi%20icin%20Malzeme.pdf}
\end{itemize}
No complete K-12 curriculum taught in T\.{I}D was found.

\section{Linguistic and Reference Materials}
\begin{itemize}
\item Ministry of Family and Social Services news item about its video-based Turkish Sign Language dictionary. The dictionary has about 2,000 common words and 11,428 videos (sign, meaning, example): \url{https://aile.gov.tr/haberler/bakanligin-guncel-turk-isaret-dili-sozlugune-150-ulkeden-4-milyondan-fazla-erisim-saglandi/}
\item Preface of the Ministry of National Education's Turkish Sign Language dictionary (Turkish): \url{https://Sozluk.Gov.Tr/dosyalar/onsoz.pdf}
\item Overview page ``\.{I}\c{s}aret Dili Nedir?'' from a Ministry of National Education school site: \url{https://niluferieml.meb.k12.tr/icerikler/isaret-dili_46306.html}
\end{itemize}

\section{Religious and Community Resources}
None were found in this search.

\section{Digital and Audio Resources}
\begin{itemize}
\item The video dictionary above is the main digital resource found. News item on the 23 September International Day of Sign Languages panel on T\.{I}D: \url{https://www.aile.gov.tr/eyhgm/haberler/23-eylul-dunya-isaret-dilleri-gunu-nde-turk-isaret-dilinin-dunu-bugunu-yarini-paneli-isitme-engellilerin-sesi-oldu/}
\end{itemize}

\section{Curriculum Development Context}
A video dictionary and some research materials exist, but no full curriculum taught through T\.{I}D was found. Further searching is needed on deaf schools, MEB programmes and the deaf federation.
\end{document}